% Adapted from Dynamic Programming by Pascal Michaillat. CC BY 4.0.
% See cover.tex for sources and adaptation details.
\documentclass{article}
\usepackage{hyperref}
\begin{document}
\newcounter{tpequation}
\setcounter{section}{3}
\section{Theory of deterministic problems}\label{sec:theory}\wcSpan*[id=ref-sec-theory]{}

So far, we have only applied dynamic programming to specific problems. In this section we propose a general treatment of dynamic programming for deterministic problems. The goal is to show you the type of problems that can be solved with dynamic programming.

Consider the following problem: Given initial condition $a_{0}$, choose $\left\lbrace c_{t}\right\rbrace _{t=0}^{\infty}$ to maximize

\[
\sum_{t=0}^{\infty}\beta^{t} u(a_{t},c_{t})
\]
subject to the law of motion

\[
a_{t+1}=g(a_{t},c_{t}).
\]
$u(a_{t},c_{t})$ is a concave function. 

Dynamic programming seeks a time-invariant policy function $h$ mapping the state $a_{t}$ into the control $c_{t}$ such that the sequence
$\lbrace c_{t}\rbrace _{t=0}^{\infty}$ generated by iterating the two functions

\[
\begin{align*}
c_{t} & =h(a_{t})\\
a_{t+1} & =g(a_{t},c_{t})
\end{align*}
\]
solves the original problem.

\subsection{Step 1: Bellman equation}

The Principle of Optimality allows us to write the value function $V$ as the solution of a functional equation:
\wcSpan*[id=equation-33]{}
\setcounter{tpequation}{32}\refstepcounter{tpequation}\label{e:bellmansymb}
\[
V(a)=\max_{c}\left(u(a,c)+\beta V(g(a,c))\right).\tag{33}
\]
This functional equation is the Bellman equation.\wcTpRef*[href=^principle-of-optimality]{} 

The optimal consumption is given by the policy function: $c=h(a)$. Another representation of the Bellman equation is 
\wcSpan*[id=equation-34]{}
\setcounter{tpequation}{33}\refstepcounter{tpequation}\label{e:bellmansymb2}
\[
V(a)=u(a,h(a))+\beta V(g(a,h(a))).\tag{34}
\]

To highlight the recursive structure of the problem, we can write the symbolic representation of the Bellman equation:

\[
V(\text{state(t)}) = \max_{\text{control(t)}}\left(u(\text{control(t)},\text{state(t)})+\beta V(\text{state(t+1)})\right)
\]
subject to

\[
\text{state(t+1)} = g(\text{control(t)},\text{state(t)}),
\]
which is equivalent to

\[
V(\text{state(t)}) = \max_{\text{control(t)}}\left(u(\text{control(t)},\text{state(t)})+\beta V(g(\text{control(t)},\text{state(t)}))\right),
\]
where control(t) and state(t) are vectors of control variables and state variables.

\subsection{Step 2: First-order condition}

Taking the first-order condition with respect to $c$ of the optimization problem~(\ref{e:bellmansymb}) yields 
\wcSpan*[id=equation-35]{}
\setcounter{tpequation}{34}\refstepcounter{tpequation}\label{e:focsymb}
\[
\frac{\partial u}{\partial c}(a,h(a))+\beta \cdot \frac{\partial g}{\partial c}(a,h(a)) \cdot \frac{d V}{d a}(g(a,h(a)))=0,\tag{35}
\]
where

\[
\frac{\partial g}{\partial c}(a,h(a))
\]
designates the partial derivative of the function $g(a,c)$ with respect to the second variable, $c$, evaluated at the pair $(a,h(a))$;


\[
\frac{\partial u}{\partial c}(a,h(a))
\]
designates the partial derivative of the function $u(a,c)$ with respect to the second variable, $c$, evaluated at the pair $(a,h(a))$;
and 

\[
\frac{d V}{d a}(g(a,h(a)))
\]
designates the derivative of the function $V(a)$ with respect to the variable $a$ evaluated at $g(a,h(a))$.

\subsection{Step 3: Benveniste-Scheinkman equation}

In the first-order condition~(\ref{e:focsymb}), we do not know the derivative $d V/d a$ of the value function (because we do not know the value function). Hence, the next step is to determine what the derivative $d V/d a$ of the value function is. To do so, we apply the Benveniste-Scheinkman theorem. This theorem says that under some regularity conditions,
\wcSpan*[id=equation-36]{}
\setcounter{tpequation}{35}\refstepcounter{tpequation}\label{e:BSET}
\[
\begin{align*}
\frac{d V}{d a}(a)=\frac{\partial u}{\partial a}(a,h(a))+\beta\cdot \frac{\partial g}{\partial a}(a,h(a))\cdot \frac{d V}{d a}(g(a,h(a))).\tag{36}
\end{align*}
\]
The theorem is a version of the envelope theorem applied to the Bellman equation~(\ref{e:bellmansymb}).

\subsection{Step 3': A combination}

The first-order condition~(\ref{e:focsymb}) yields
\wcSpan*[id=equation-37]{}

\[
\begin{align*}
\frac{d V}{d a}(g(a,h(a)))=-\frac{1}{\beta}\cdot \frac{\partial u(a,h(a))/\partial c}{\partial g(a,h(a))/\partial c}.\tag{37}
\end{align*}
\]
Combining this equation with the Benveniste-Scheinkman equation~(\ref{e:BSET}) yields
\wcSpan*[id=equation-38]{}
\setcounter{tpequation}{37}\refstepcounter{tpequation}\label{e:plug}
\[
\begin{align*}
\frac{d V}{d a}(a)=\frac{\partial u}{\partial a}(a,h(a))-\frac{\partial u}{\partial c}(a,h(a))\cdot \frac{\partial g(a,h(a))/\partial a}{\partial g(a,h(a))/\partial c}.\tag{38}
\end{align*}
\]
This step is necessary when 

\[
\frac{\partial g(a,c)}{\partial a}\neq 0.
\] In the consumption-saving problem, we picked a control variable---next period's wealth, $a'$---such that the state variable, $a$, does not enter the transition function, $g$. Therefore, we could bypass this step.

\subsection{Step 4: One step forward}

Equation (\ref{e:plug}) is true for any value of the state variable $a$. In particular, it is true for $a'=g(a,h(a))$. Therefore
\wcSpan*[id=equation-39]{}
\setcounter{tpequation}{38}\refstepcounter{tpequation}\label{e:onestep}
\[
\frac{d V}{d a}(a')=\frac{\partial u}{\partial a}(a',h(a'))-\frac{\partial u}{\partial c}(a',h(a'))\cdot \frac{\partial g(a',h(a'))/\partial a}{\partial g(a',h(a'))/\partial c}.\tag{39}
\]

\subsection{Step 5: Euler equation}

We plug equation~(\ref{e:onestep}) into equation~(\ref{e:focsymb}) to get the Euler equation

\[
\frac{\partial u}{\partial c}(a,h(a))+\beta \cdot \frac{\partial g}{\partial c}(a,h(a)) \left[\frac{\partial u}{\partial a}(a',h(a'))-\frac{\partial u}{\partial c}(a',h(a')) \cdot \frac{\partial g(a',h(a'))/\partial a}{\partial g(a',h(a'))/\partial c}\right]=0.
\]
The equation characterizes the optimal behavior of the control variable $c$.


\wcTpNote[ref=principle-of-optimality,title=Principle of optimality]{The proof of the Principle of Optimality is due to Bellman. The formal derivation and proof of this result, as well as the conditions under which this result holds, are omitted here.}
\subsection*{Note}
\wcTpListof[selector=tp-note]{}

\end{document}
